\documentclass[12pt]{article}
\usepackage[spanish]{babel}
\usepackage[T1]{fontenc}
\usepackage[utf8]{inputenc}
\usepackage{amsmath,amssymb}
\usepackage{amsthm}
\usepackage{geometry}
\geometry{margin=2.5cm}
\usepackage{comment}
\renewcommand{\familydefault}{\sfdefault} % texto en sans serif
% \usepackage{lmodern} % alternativa si prefieres serif moderna
\usepackage{mathrsfs}

\newtheorem{definicion}{Definición}[section]
\newtheorem{proposicion}{Proposición}[section]
\newtheorem{teorema}{Teorema}[section]
\newtheorem{corolario}{Corolario}[section]
\newtheorem{observacion}{Observación}[section]
\newtheorem{ejemplo}{Ejemplo}[section]

% Figuras y tablas
\usepackage{graphicx}
\usepackage{float}       % para [H]
\usepackage{subcaption}
\usepackage{booktabs,tabularx,array}

% Listas
\usepackage{enumitem}

% Títulos (opcional; quítalo si no lo usas o si cambias a wiley-article)
\usepackage{titlesec}

% Bibliografía con natbib (autor-año)
\usepackage[authoryear]{natbib}
\bibliographystyle{plainnat}

% Hipervínculos (SIEMPRE al final)
\usepackage[hidelinks]{hyperref}
\title{Control Informacional}

\author{
\textbf{Alejandro Cruz López} \\
Universidad Autónoma Metropolitana \\
Unidad Iztapalapa
}

\date{
\vspace{1em}
\textbf{Asesores} \\[0.3em]
Dr.\ Asael Fabián Martínez Martínez \\
Dr.\ Joaquín Delgado Fernández \\
\vspace{1em}
9 de abril de 2026
}
\begin{document}
\maketitle
\section{Markov y control informacional}
\begin{definicion}[Matriz estocástica y operador de Markov]
Sea \(n\geq 2\). Una matriz \(P\in\mathbb{R}^{n\times n}\) se dice
\emph{estocástica por filas} si
\[
P_{ij}\geq 0
\quad\text{para todo } i,j\in\{1,\dots,n\},
\]
y
\[
\sum_{j=1}^{n}P_{ij}=1
\quad\text{para todo } i\in\{1,\dots,n\}.
\]

Dada una matriz estocástica por filas \(P\), se define el operador de Markov
asociado a \(P\) como la aplicación
\[
L_P:\Delta_n\longrightarrow \Delta_n,
\qquad
L_P(X):=XP,
\]
donde \(X\in\Delta_n\) se interpreta como vector fila. En coordenadas,
\[
\bigl(L_P(X)\bigr)_j
=
\sum_{i=1}^{n}X_iP_{ij},
\qquad
j=1,\dots,n.
\]
\end{definicion}
\begin{ejemplo}[Operador de Markov en dimensión \(3\)]
Considérese el estado
\[
X=\left(\frac{1}{2},\frac{1}{3},\frac{1}{6}\right)\in\Delta_3
\]
y la matriz estocástica por filas
\[
P=
\begin{pmatrix}
\frac{1}{2} & \frac{1}{2} & 0 \\[0.3em]
\frac{1}{4} & \frac{1}{2} & \frac{1}{4} \\[0.3em]
0 & \frac{1}{3} & \frac{2}{3}
\end{pmatrix}.
\]
En efecto, todas las entradas de \(P\) son no negativas y cada fila suma \(1\).

El operador de Markov asociado está dado por
\[
L_P(X)=XP.
\]
Calculando coordenada por coordenada,
\[
\bigl(L_P(X)\bigr)_1
=
\frac{1}{2}\cdot\frac{1}{2}
+
\frac{1}{3}\cdot\frac{1}{4}
+
\frac{1}{6}\cdot 0
=
\frac{1}{4}+\frac{1}{12}
=
\frac{1}{3},
\]
\[
\bigl(L_P(X)\bigr)_2
=
\frac{1}{2}\cdot\frac{1}{2}
+
\frac{1}{3}\cdot\frac{1}{2}
+
\frac{1}{6}\cdot\frac{1}{3}
=
\frac{1}{4}+\frac{1}{6}+\frac{1}{18}
=
\frac{17}{36},
\]
y
\[
\bigl(L_P(X)\bigr)_3
=
\frac{1}{2}\cdot 0
+
\frac{1}{3}\cdot\frac{1}{4}
+
\frac{1}{6}\cdot\frac{2}{3}
=
\frac{1}{12}+\frac{1}{9}
=
\frac{7}{36}.
\]
Por tanto,
\[
L_P(X)
=
\left(
\frac{1}{3},
\frac{17}{36},
\frac{7}{36}
\right).
\]
Finalmente,
\[
\frac{1}{3}+\frac{17}{36}+\frac{7}{36}
=
\frac{12}{36}+\frac{17}{36}+\frac{7}{36}
=
1,
\]
de modo que \(L_P(X)\in\Delta_3\).
\end{ejemplo}

\begin{proposicion}[Preservación del símplex]
Sea \(P\in\mathbb{R}^{n\times n}\) una matriz estocástica por filas. Entonces:

\begin{enumerate}[label=\textnormal{(\roman*)}]
    \item Para todo \(X\in\Delta_n\), se cumple
    \[
    L_P(X)\in\Delta_n.
    \]

    \item Si \(X\in\Delta_n^\circ\) y \(P_{ij}>0\) para todo \(i,j\), entonces
    \[
    L_P(X)\in\Delta_n^\circ.
    \]

    \item El operador \(L_P\) es afín; es decir, para todo
    \(X,Z\in\Delta_n\) y todo \(\alpha\in[0,1]\),
    \[
    L_P\bigl(\alpha X+(1-\alpha)Z\bigr)
    =
    \alpha L_P(X)+(1-\alpha)L_P(Z).
    \]
\end{enumerate}
\end{proposicion}

\begin{proof}
Sea \(X\in\Delta_n\). Para cada \(j\in\{1,\dots,n\}\),
\[
\bigl(L_P(X)\bigr)_j
=
\sum_{i=1}^{n}X_iP_{ij}.
\]
Como \(X_i\geq 0\) y \(P_{ij}\geq 0\), se tiene
\[
\bigl(L_P(X)\bigr)_j\geq 0.
\]
Además,
\[
\sum_{j=1}^{n}\bigl(L_P(X)\bigr)_j
=
\sum_{j=1}^{n}\sum_{i=1}^{n}X_iP_{ij}
=
\sum_{i=1}^{n}X_i\sum_{j=1}^{n}P_{ij}.
\]
Como \(P\) es estocástica por filas,
\[
\sum_{j=1}^{n}P_{ij}=1
\]
para todo \(i\). Por tanto,
\[
\sum_{j=1}^{n}\bigl(L_P(X)\bigr)_j
=
\sum_{i=1}^{n}X_i
=
1.
\]
Así, \(L_P(X)\in\Delta_n\).

Ahora supóngase que \(X\in\Delta_n^\circ\) y que \(P_{ij}>0\) para todo \(i,j\).
Entonces \(X_i>0\) para todo \(i\), y por tanto
\[
\bigl(L_P(X)\bigr)_j
=
\sum_{i=1}^{n}X_iP_{ij}
>
0
\]
para cada \(j\). En consecuencia,
\[
L_P(X)\in\Delta_n^\circ.
\]

Finalmente, sean \(X,Z\in\Delta_n\) y \(\alpha\in[0,1]\). Por linealidad del producto
matricial,
\[
L_P\bigl(\alpha X+(1-\alpha)Z\bigr)
=
\bigl(\alpha X+(1-\alpha)Z\bigr)P
\]
\[
=
\alpha XP+(1-\alpha)ZP
=
\alpha L_P(X)+(1-\alpha)L_P(Z).
\]
Esto prueba que \(L_P\) es afín.
\end{proof}

\begin{definicion}[Corrección local inducida por Markov]
Sea \(P\in\mathbb{R}^{n\times n}\) una matriz estocástica por filas y sea
\(X\in\Delta_n^\circ\). Se define la \emph{corrección local inducida por Markov}
en el estado \(X\) como el vector
\[
Y^P(X)\in\mathbb{R}^n
\]
cuyas coordenadas están dadas por
\[
Y_i^P(X)
:=
\frac{(L_P(X))_i}{X_i}-1,
\qquad i=1,\dots,n.
\]
Equivalentemente, como \(L_P(X)=XP\),
\[
Y_i^P(X)
=
\frac{(XP)_i}{X_i}-1
=
\frac{\sum_{k=1}^{n}X_kP_{ki}}{X_i}-1.
\]
\end{definicion}

\begin{proposicion}[Compatibilidad local entre Markov y CI]
Sea \(P\in\mathbb{R}^{n\times n}\) una matriz estocástica por filas y sea
\(X\in\Delta_n^\circ\). Entonces la corrección local inducida por Markov satisface
\[
Y^P(X)\in C(X).
\]
Además,
\[
T^{\mathrm{loc}}_{Y^P(X)}(X)=L_P(X).
\]
\end{proposicion}

\begin{proof}
Por definición,
\[
Y_i^P(X)
=
\frac{(L_P(X))_i}{X_i}-1.
\]
Como \(X\in\Delta_n^\circ\), se tiene \(X_i>0\) para todo \(i\), de modo que la expresión
anterior está bien definida.

Primero verificamos la admisibilidad local. Para cada \(i\),
\[
X_i\bigl(1+Y_i^P(X)\bigr)
=
X_i\left(1+\frac{(L_P(X))_i}{X_i}-1\right)
=
(L_P(X))_i.
\]
Por la preservación del símplex bajo \(L_P\), se tiene
\[
(L_P(X))_i\geq 0
\]
para todo \(i\). Por tanto,
\[
X_i\bigl(1+Y_i^P(X)\bigr)\geq 0.
\]
Además,
\[
\sum_{i=1}^{n}X_i\bigl(1+Y_i^P(X)\bigr)
=
\sum_{i=1}^{n}(L_P(X))_i
=
1.
\]
Así,
\[
Y^P(X)\in C(X).
\]

Finalmente, por la definición del operador local de CI,
\[
\bigl(T^{\mathrm{loc}}_{Y^P(X)}(X)\bigr)_i
=
X_i\bigl(1+Y_i^P(X)\bigr).
\]
Usando el cálculo anterior,
\[
\bigl(T^{\mathrm{loc}}_{Y^P(X)}(X)\bigr)_i
=
(L_P(X))_i.
\]
Como esto vale para todo \(i\), se concluye que
\[
T^{\mathrm{loc}}_{Y^P(X)}(X)=L_P(X).
\]
\end{proof}
\begin{observacion}[Dependencia local de la corrección markoviana]
La corrección \(Y^P(X)\) depende, en general, del estado base \(X\). En efecto,
\[
1+Y_i^P(X)
=
\frac{(XP)_i}{X_i}.
\]
Por tanto, para que existiera un vector fijo \(Y\in C_{\mathrm{glob}}\), independiente de
\(X\), tal que
\[
T_Y(X)=L_P(X)
\qquad
\text{para todo } X\in\Delta_n^\circ,
\]
sería necesario que los cocientes
\[
\frac{(XP)_i}{X_i}
\]
fueran independientes de \(X\), salvo por el factor común de normalización propio de
\(T_Y\). Esto no ocurre en general, pues \((XP)_i\) mezcla las coordenadas de \(X\) mediante
la columna \(i\) de \(P\), mientras que \(X_i\) sólo depende de la coordenada \(i\).

En consecuencia, un operador de Markov \(L_P\) debe interpretarse, dentro de CI, como una
familia de correcciones locales
\[
X\longmapsto Y^P(X),
\]
no como una única corrección global fija del grupo \(G\), salvo en casos particulares.
La única excepción ocurre cuando $P$ tiene todas las filas iguales, es decir,
$P_{ij}=p_j$ para todo $i$ (matriz de rango~$1$). En ese caso
$(XP)_i = p_i$ para todo $X$, de modo que $1+Y_i^P(X)=p_i/X_i$ depende
de $X$ pero el operador $T_Y$ con $Y_i=p_i/X_i-1$ envía cualquier $X$
al mismo punto $p\in\Delta_n^\circ$. Para toda otra $P$, no existe
$Y\in C_{\mathrm{glob}}$ constante tal que $T_Y(X)=XP$ para todo
$X\in\Delta_n^\circ$.

\end{observacion}




\begin{definicion}[Distribución estacionaria de una matriz de Markov]
Sea \(P\in\mathbb{R}^{n\times n}\) una matriz estocástica por filas. 
Una distribución \(\pi\in\Delta_n\) se dice \emph{estacionaria} para \(P\) si
\[
\pi P=\pi.
\]

Equivalente, \(\pi\) es estacionaria si es un punto fijo del operador de Markov
\[
L_P:\Delta_n\longrightarrow\Delta_n,
\qquad
L_P(X)=XP.
\]
Es decir,
\[
L_P(\pi)=\pi.
\]

En coordenadas, la condición de estacionariedad se escribe como
\[
\pi_j
=
\sum_{i=1}^{n}\pi_iP_{ij},
\qquad
j=1,\dots,n.
\]
\end{definicion}

\begin{proposicion}[Estacionariedad y corrección nula]
Sea \(P\in\mathbb{R}^{n\times n}\) una matriz estocástica por filas y sea
\(\pi\in\Delta_n^\circ\) una distribución estacionaria de \(P\). Entonces
\[
Y^P(\pi)=0.
\]
En consecuencia,
\[
T^{\mathrm{loc}}_{Y^P(\pi)}(\pi)=\pi.
\]
\end{proposicion}

\begin{proof}
Como $\pi$ es estacionaria, $(\pi P)_i = \pi_i$ para todo $i$. Por definición,
\[
Y_i^P(\pi)
=
\frac{(\pi P)_i}{\pi_i}-1
=
\frac{\pi_i}{\pi_i}-1
=
0.
\]
Como esto vale para todo $i$, se concluye $Y^P(\pi)=0$.
La segunda afirmación se sigue de la compatibilidad local:
$T^{\mathrm{loc}}_{Y^P(\pi)}(\pi)=L_P(\pi)=\pi P=\pi$.
\end{proof}

\begin{proposicion}[Trayectoria markoviana como trayectoria local de CI]
Sea \(P\in\mathbb{R}^{n\times n}\) una matriz estocástica por filas y sea
\(X_0\in\Delta_n^\circ\). Defínase la trayectoria markoviana
\[
X_{t+1}=L_P(X_t)=X_tP,
\qquad t=0,1,2,\dots.
\]
Entonces cada paso de la trayectoria puede escribirse como una corrección local de CI:
\[
X_{t+1}
=
T^{\mathrm{loc}}_{Y^P(X_t)}(X_t).
\]
En particular,
\[
X_t=L_P^t(X_0)=X_0P^t.
\]
\end{proposicion}

\begin{proof}
Por definición de la trayectoria,
\[
X_{t+1}=L_P(X_t).
\]
Como \(X_t\in\Delta_n^\circ\), la corrección local inducida por Markov está bien definida:
\[
Y_i^P(X_t)
=
\frac{(X_tP)_i}{(X_t)_i}-1.
\]
Por la compatibilidad local entre \(L_P\) y \(T^{\mathrm{loc}}\), se tiene
\[
T^{\mathrm{loc}}_{Y^P(X_t)}(X_t)=L_P(X_t).
\]
Por tanto,
\[
X_{t+1}
=
T^{\mathrm{loc}}_{Y^P(X_t)}(X_t).
\]
La identidad
\[
X_t=X_0P^t
\]
se obtiene por inducción. Para \(t=0\) es inmediata. Si \(X_t=X_0P^t\), entonces
\[
X_{t+1}=X_tP=(X_0P^t)P=X_0P^{t+1}.
\]
\end{proof}

\begin{proposicion}[La matriz fija genera un campo vectorial discreto]
Sea \(P\in\mathbb{R}^{n\times n}\) una matriz estocástica por filas. El operador
\(L_P\) induce sobre el símplex el campo discreto
\[
V_P(X):=L_P(X)-X=XP-X.
\]
Entonces
\[
\sum_{i=1}^n (V_P(X))_i=0.
\]
Por tanto,
\[
V_P(X)\in T_X\Delta_n^\circ,
\]
es decir, el campo \(V_P\) es tangente al símplex.
\end{proposicion}

\begin{proof}
Como \(X\in\Delta_n\) y \(P\) es estocástica por filas, se tiene
\[
L_P(X)=XP\in\Delta_n.
\]
Por tanto,
\[
\sum_{i=1}^n (XP)_i=1
\qquad\text{y}\qquad
\sum_{i=1}^n X_i=1.
\]
Entonces
\[
\sum_{i=1}^n (V_P(X))_i
=
\sum_{i=1}^n \bigl((XP)_i-X_i\bigr)
=
\sum_{i=1}^n (XP)_i-\sum_{i=1}^n X_i
=
1-1
=
0.
\]
Así, \(V_P(X)\) pertenece al hiperplano tangente del símplex.
\end{proof}

\begin{proposicion}[Disipación de energía relativa hacia la estacionaria]
Sea \(P\in\mathbb{R}^{n\times n}\) una matriz estocástica por filas y sea
\(\pi\in\Delta_n^\circ\) una distribución estacionaria de \(P\). Defínase
\[
E_\pi(X):=D_{\mathrm{KL}}(\pi\|X)
=
\sum_{i=1}^n \pi_i\log\left(\frac{\pi_i}{X_i}\right).
\]
Entonces, para todo \(X\in\Delta_n^\circ\),
\[
E_\pi(L_P(X))\leq E_\pi(X).
\]
Equivalentemente,
\[
D_{\mathrm{KL}}(\pi\|XP)
\leq
D_{\mathrm{KL}}(\pi\|X).
\]
\end{proposicion}
\begin{proof}
Como $\pi$ es estacionaria, $\pi_j=\sum_i \pi_i P_{ij}$ para todo $j$. Por tanto,
\[
D_{\mathrm{KL}}(\pi\|XP)
=
\sum_j \pi_j \log\frac{\pi_j}{(XP)_j}
=
\sum_j \left(\sum_i \pi_i P_{ij}\right)
\log\frac{\sum_i \pi_i P_{ij}}{\sum_i X_i P_{ij}}.
\]
Para cada $j$ fijo, los pesos $w_{ij}:=\pi_i P_{ij}/\pi_j$ satisfacen
$w_{ij}\geq 0$ y $\sum_i w_{ij}=1$. Aplicando Jensen a la función cóncava $\log$,
\[
\log\frac{\sum_i \pi_i P_{ij}}{\sum_i X_i P_{ij}}
\leq
\sum_i w_{ij}\log\frac{\pi_i}{X_i}.
\]
Multiplicando por $\pi_j$ y sumando en $j$,
\[
D_{\mathrm{KL}}(\pi\|XP)
\leq
\sum_j\sum_i \pi_i P_{ij}\log\frac{\pi_i}{X_i}
=
\sum_i \pi_i\log\frac{\pi_i}{X_i}
\underbrace{\sum_j P_{ij}}_{=1}
=
D_{\mathrm{KL}}(\pi\|X).
\]
Es decir, $E_\pi(L_P(X))\leq E_\pi(X)$.
\end{proof}

\begin{proposicion}[Composición de operadores de Markov]
Sean \(P,Q\in\mathbb{R}^{n\times n}\) matrices estocásticas por filas. Entonces \(PQ\)
también es una matriz estocástica por filas y
\[
L_Q\circ L_P=L_{PQ}.
\]
Es decir, para todo \(X\in\Delta_n\),
\[
L_Q(L_P(X))=L_{PQ}(X).
\]
\end{proposicion}

\begin{proof}
Primero verificamos que \(PQ\) es estocástica por filas. Como \(P_{ik}\geq0\) y
\(Q_{kj}\geq0\), se tiene
\[
(PQ)_{ij}=\sum_{k=1}^n P_{ik}Q_{kj}\geq 0.
\]
Además,
\[
\sum_{j=1}^n (PQ)_{ij}
=
\sum_{j=1}^n\sum_{k=1}^n P_{ik}Q_{kj}
=
\sum_{k=1}^n P_{ik}\sum_{j=1}^n Q_{kj}.
\]
Como \(Q\) es estocástica por filas,
\[
\sum_{j=1}^n Q_{kj}=1.
\]
Así,
\[
\sum_{j=1}^n (PQ)_{ij}
=
\sum_{k=1}^n P_{ik}
=
1.
\]
Por tanto, \(PQ\) es estocástica por filas.

Ahora, para todo \(X\in\Delta_n\),
\[
L_Q(L_P(X))
=
L_Q(XP)
=
(XP)Q
=
X(PQ)
=
L_{PQ}(X).
\]
\end{proof}


\subsection{Dinámica de matrices estocásticas}

En esta sección se formaliza únicamente lo que fue explorado en los dos
experimentos computacionales anteriores: primero, una dinámica de relajación de
una matriz estocástica hacia una matriz objetivo; segundo, una dinámica extrema
en la que la matriz cambia bajo un campo de velocidades y se proyecta de nuevo
al conjunto de matrices estocásticas. No se introduce todavía ningún componente adicional ajeno a estos dos
experimentos.

\begin{definicion}[Dinámica markoviana no autónoma]
Sea \(\mathcal S_n\) el conjunto de matrices estocásticas por filas de tamaño
\(n\times n\). Una dinámica markoviana no autónoma sobre \(\Delta_n\) es una
sucesión
\[
(X_t,K_t)_{t\geq 0}
\]
tal que
\[
X_t\in\Delta_n,
\qquad
K_t\in\mathcal S_n,
\]
y
\[
X_{t+1}=X_tK_t,
\qquad t\geq 0.
\]
El vector \(X_t\) representa el estado informacional observado en el tiempo
\(t\), mientras que \(K_t\) representa el mecanismo estocástico activo en ese
mismo tiempo.
\end{definicion}

\begin{proposicion}[Cierre del estado bajo dinámica no autónoma]
Si \(X_0\in\Delta_n\) y \(K_t\in\mathcal S_n\) para todo \(t\geq 0\), entonces
\[
X_t\in\Delta_n
\qquad
\text{para todo }t\geq 0.
\]
\end{proposicion}

\begin{proof}
Supongamos que \(X_t\in\Delta_n\). Como \(K_t\in\mathcal S_n\), se tiene
\((K_t)_{ij}\geq 0\). Por tanto,
\[
(X_tK_t)_j=\sum_{i=1}^n X_{t,i}(K_t)_{ij}\geq 0.
\]
Además,
\[
\sum_{j=1}^n (X_tK_t)_j
=
\sum_{j=1}^n\sum_{i=1}^n X_{t,i}(K_t)_{ij}
=
\sum_{i=1}^n X_{t,i}\sum_{j=1}^n (K_t)_{ij}
=
\sum_{i=1}^n X_{t,i}
=
1.
\]
Así, \(X_{t+1}=X_tK_t\in\Delta_n\). La conclusión se obtiene por inducción.
\end{proof}

\begin{proposicion}[Producto temporal ordenado]
Sea \((X_t,K_t)_{t\geq 0}\) una dinámica markoviana no autónoma. Entonces,
para todo \(t\geq 1\),
\[
X_t=X_0K_0K_1\cdots K_{t-1}.
\]
\end{proposicion}

\begin{proof}
Para \(t=1\), la identidad es \(X_1=X_0K_0\). Supongamos que
\[
X_t=X_0K_0K_1\cdots K_{t-1}.
\]
Entonces,
\[
X_{t+1}=X_tK_t
=
X_0K_0K_1\cdots K_{t-1}K_t.
\]
Por inducción, la fórmula vale para todo \(t\geq 1\).
\end{proof}

\begin{definicion}[Matriz estocástica primitiva]
Una matriz \(K\in\mathcal S_n\) se dice \emph{primitiva} si existe
\(m\geq 1\) tal que
\[
(K^m)_{ij}>0
\qquad
\text{para todo }i,j\in\{1,\dots,n\}.
\]
En particular, toda matriz primitiva es irreducible y admite una única
distribución estacionaria \(\pi\in\Delta_n^\circ\).
\end{definicion}



\begin{definicion}[Estacionaria instantánea]
Sea \(K_t\in\mathcal S_n\). Una distribución \(\pi_t\in\Delta_n\) se llama
estacionaria instantánea de \(K_t\) si
\[
\pi_tK_t=\pi_t.
\]
Cuando \(K_t\) es primitiva, esta distribución es única y pertenece a
\(\Delta_n^\circ\). En ese caso escribiremos
\[
\pi_t=\pi(K_t).
\]
\end{definicion}

\begin{definicion}[Observables del desfase dinámico]
Supongamos que cada \(K_t\) admite una estacionaria instantánea única
\(\pi_t\in\Delta_n^\circ\). Definimos:
\[
\mathcal E_t
:=
D_{\mathrm{KL}}(\pi_t\Vert X_t)
=
\sum_{i=1}^n \pi_{t,i}\log\frac{\pi_{t,i}}{X_{t,i}},
\]
\[
\ell_t:=\|X_t-\pi_t\|_1,
\]
\[
\nu_t:=\|K_t-K_{t-1}\|_F,
\qquad t\geq 1,
\]
\[
\chi_t:=\|K_tK_{t-1}-K_{t-1}K_t\|_F,
\qquad t\geq 1.
\]
Aquí \(\mathcal E_t\) mide desfase informacional, \(\ell_t\) mide lag en norma
\(1\), \(\nu_t\) mide velocidad matricial y \(\chi_t\) mide no conmutatividad
temporal.
\end{definicion}

\begin{proposicion}[Control del lag por divergencia KL]
Para cada \(t\), si \(X_t\in\Delta_n^\circ\) y \(\pi_t\in\Delta_n^\circ\),
entonces
\[
\ell_t^2
\leq
2\mathcal E_t.
\]
En particular, si \(\mathcal E_t\leq\varepsilon\), entonces
\[
\ell_t\leq \sqrt{2\varepsilon}.
\]
\end{proposicion}

\begin{proof}
La desigualdad es una aplicación directa de Pinsker:
\[
D_{\mathrm{KL}}(\pi_t\Vert X_t)
\geq
\frac12\|\pi_t-X_t\|_1^2.
\]
Sustituyendo las definiciones de \(\mathcal E_t\) y \(\ell_t\), se obtiene
\[
\mathcal E_t\geq \frac12\ell_t^2.
\]
Por tanto, \(\ell_t^2\leq 2\mathcal E_t\), y la segunda afirmación se sigue
tomando raíz cuadrada.
\end{proof}

\begin{proposicion}[La energía instantánea no es necesariamente monótona]
Sea $\pi\in\Delta_n^\circ$ fija. En una dinámica autónoma $X_{t+1}=X_tP$
con $P$ fija y $\pi$ su estacionaria, la sucesión $E_\pi(X_t)=D_{\mathrm{KL}}(\pi\|X_t)$
es monótona decreciente (Proposición~anterior). Sin embargo, en una dinámica
no autónoma, la sucesión
\[
\mathcal E_t:=D_{\mathrm{KL}}(\pi_t\Vert X_t),
\]
donde $\pi_t=\pi(K_t)$ es la estacionaria \emph{instantánea} de $K_t$,
no tiene por qué ser monótona decreciente. La no monotonicidad no refleja
un crecimiento intrínseco de la divergencia, sino el hecho de que el
objetivo $\pi_t$ se desplaza.
\end{proposicion}

\begin{proof}
Basta exhibir un contraejemplo. Tomemos \(n=2\) y
\[
X_0=\left(\frac12,\frac12\right).
\]
Sean
\[
K_0=
\begin{pmatrix}
0.9 & 0.1\\
0.9 & 0.1
\end{pmatrix},
\qquad
K_1=
\begin{pmatrix}
0.1 & 0.9\\
0.1 & 0.9
\end{pmatrix}.
\]
Entonces las estacionarias instantáneas son
\[
\pi_0=(0.9,0.1),
\qquad
\pi_1=(0.1,0.9).
\]
Además,
\[
X_1=X_0K_0=(0.9,0.1)=\pi_0.
\]
Por tanto,
\[
\mathcal E_0
=
D_{\mathrm{KL}}\left((0.9,0.1)\middle\Vert
\left(\frac12,\frac12\right)\right)
\approx 0.368,
\]
mientras que
\[
\mathcal E_1
=
D_{\mathrm{KL}}\bigl((0.1,0.9)\Vert(0.9,0.1)\bigr)
\approx 1.758.
\]
Así,
\[
\mathcal E_1>\mathcal E_0.
\]
Luego la energía instantánea puede aumentar cuando cambia la matriz de
transición.
\end{proof}

\subsection{Primer experimento: relajación hacia una matriz objetivo}

\begin{definicion}[Relajación matricial estocástica]
Sean \(K_0,K_\ast\in\mathcal S_n\) y sea \(\eta\in(0,1)\). La relajación de
\(K_0\) hacia \(K_\ast\) se define por
\[
K_{t+1}=(1-\eta)K_t+\eta K_\ast,
\qquad t\geq 0.
\]
Acoplada a esta evolución matricial, se considera la dinámica de estados
\[
X_{t+1}=X_tK_t.
\]
\end{definicion}

\begin{proposicion}[Cierre estocástico de la relajación]
Si \(K_0,K_\ast\in\mathcal S_n\), entonces
\[
K_t\in\mathcal S_n
\qquad
\text{para todo }t\geq 0.
\]
\end{proposicion}

\begin{proof}
Supongamos que \(K_t\in\mathcal S_n\). Para cada \(i,j\),
\[
(K_{t+1})_{ij}
=
(1-\eta)(K_t)_{ij}+\eta(K_\ast)_{ij}
\geq 0.
\]
Además, para cada fila \(i\),
\[
\sum_{j=1}^n (K_{t+1})_{ij}
=
(1-\eta)\sum_{j=1}^n(K_t)_{ij}
+
\eta\sum_{j=1}^n(K_\ast)_{ij}
=
(1-\eta)+\eta
=
1.
\]
Por tanto \(K_{t+1}\in\mathcal S_n\). Como \(K_0\in\mathcal S_n\), la
conclusión sigue por inducción.
\end{proof}

\begin{proposicion}[Solución explícita de la relajación]
Para todo \(t\geq 0\),
\[
K_t
=
K_\ast+(1-\eta)^t(K_0-K_\ast).
\]
\end{proposicion}

\begin{proof}
Para \(t=0\), la identidad es inmediata. Supongamos que
\[
K_t=K_\ast+(1-\eta)^t(K_0-K_\ast).
\]
Entonces,
\begin{align*}
K_{t+1}
&=(1-\eta)K_t+\eta K_\ast\\
&=(1-\eta)\bigl[K_\ast+(1-\eta)^t(K_0-K_\ast)\bigr]+\eta K_\ast\\
&=K_\ast+(1-\eta)^{t+1}(K_0-K_\ast).
\end{align*}
Por inducción, la fórmula vale para todo \(t\geq 0\).
\end{proof}

\begin{corolario}[Convergencia geométrica hacia la matriz objetivo]
En la relajación matricial estocástica,
\[
\|K_t-K_\ast\|_F
=
(1-\eta)^t\|K_0-K_\ast\|_F.
\]
En particular,
\[
\lim_{t\to\infty}K_t=K_\ast.
\]
Además, para \(t\geq 1\),
\[
\nu_t
=
\eta(1-\eta)^{t-1}\|K_0-K_\ast\|_F.
\]
\end{corolario}

\begin{proof}
Por la proposición anterior,
\[
K_t-K_\ast=(1-\eta)^t(K_0-K_\ast).
\]
Tomando norma de Frobenius,
\[
\|K_t-K_\ast\|_F
=
(1-\eta)^t\|K_0-K_\ast\|_F.
\]
Como \(0<1-\eta<1\), se obtiene \(K_t\to K_\ast\). Por otra parte,
\[
K_t-K_{t-1}
=
\bigl[K_\ast+(1-\eta)^t(K_0-K_\ast)\bigr]
-
\bigl[K_\ast+(1-\eta)^{t-1}(K_0-K_\ast)\bigr],
\]
de donde
\[
K_t-K_{t-1}
=
-\eta(1-\eta)^{t-1}(K_0-K_\ast).
\]
Tomando norma de Frobenius se obtiene la fórmula para \(\nu_t\).
\end{proof}

\begin{observacion}[Lectura del primer experimento]
El primer script no cambia arbitrariamente las probabilidades de \(X_t\).
Cambia el mecanismo \(K_t\) por interpolación convexa hacia una matriz objetivo
\(K_\ast\). La convergencia observada de \(\|K_t-K_\ast\|_F\) no es empírica:
es consecuencia directa de la fórmula cerrada
\[
K_t=K_\ast+(1-\eta)^t(K_0-K_\ast).
\]
El estado \(X_t\), en cambio, no sigue una sola matriz fija, sino el producto
ordenado
\[
X_t=X_0K_0K_1\cdots K_{t-1}.
\]
\end{observacion}

\subsection{Segundo experimento: dinámica extrema proyectada}

\begin{definicion}[Proyección fila a fila sobre el símplex]
Para \(a\in\mathbb R^n\), definimos la proyección euclidiana sobre el símplex
por
\[
\Pi_{\Delta_n}(a)
:=
\operatorname*{argmin}_{p\in\Delta_n}\|p-a\|_2.
\]
Para una matriz \(A\in\mathbb R^{n\times n}\), definimos
\(\Pi_{\mathcal S_n}(A)\) fila por fila mediante
\[
\bigl(\Pi_{\mathcal S_n}(A)\bigr)_{i\cdot}
:=
\Pi_{\Delta_n}(A_{i\cdot}),
\qquad i=1,\dots,n.
\]
\end{definicion}

\begin{proposicion}[La proyección produce matrices estocásticas]
Para toda matriz \(A\in\mathbb R^{n\times n}\),
\[
\Pi_{\mathcal S_n}(A)\in\mathcal S_n.
\]
\end{proposicion}

\begin{proof}
Por definición, cada fila de \(\Pi_{\mathcal S_n}(A)\) pertenece a
\(\Delta_n\). Por tanto, todas sus entradas son no negativas y cada fila suma
uno. Luego \(\Pi_{\mathcal S_n}(A)\in\mathcal S_n\).
\end{proof}

\begin{definicion}[Dinámica extrema proyectada]
Sea \(K_0\in\mathcal S_n\), \(X_0\in\Delta_n\), \(K_{-1}:=K_0\), y sean
\(\alpha,\beta,\mu,\rho,\tau\geq 0\). Supóngase dada una sucesión de objetivos
móviles \(M_t\in\mathcal S_n\), una retroalimentación matricial
\(F:\Delta_n\to\mathbb R^{n\times n}\) y una perturbación determinista
\(R_t\in\mathbb R^{n\times n}\). Definimos
\[
V_t
:=
\alpha(M_t-K_t)
+
\beta F(X_t)
+
\mu(K_t-K_{t-1})
+
\rho R_t,
\]
\[
K_{t+1}
:=
\Pi_{\mathcal S_n}(K_t+\tau V_t),
\]
y
\[
X_{t+1}:=X_tK_t.
\]
A esta construcción la llamaremos dinámica extrema proyectada. El término
\(M_t-K_t\) representa atracción hacia un objetivo móvil; \(F(X_t)\),
retroalimentación desde el estado; \(K_t-K_{t-1}\), inercia; y \(R_t\), una
perturbación determinista externa.
\end{definicion}

\begin{proposicion}[Cierre de la dinámica extrema]
En la dinámica extrema proyectada,
\[
K_t\in\mathcal S_n
\qquad
\text{y}
\qquad
X_t\in\Delta_n
\]
para todo \(t\geq 0\).
\end{proposicion}

\begin{proof}
Como \(K_0\in\mathcal S_n\), supongamos que \(K_t\in\mathcal S_n\). Entonces
\(K_t+\tau V_t\) es una matriz real cualquiera, pero
\[
K_{t+1}=\Pi_{\mathcal S_n}(K_t+\tau V_t)\in\mathcal S_n
\]
por la proposición anterior. Por inducción, \(K_t\in\mathcal S_n\) para todo
\(t\). Una vez probado esto, el cierre \(X_t\in\Delta_n\) se sigue de la
proposición de cierre del estado bajo dinámica no autónoma.
\end{proof}

\begin{definicion}[Presión de seguimiento]
Cuando existe \(\pi_t=\pi(K_t)\), definimos la presión de seguimiento por
\[
\omega_t:=\|\pi_{t+1}-\pi_t\|_1.
\]
Esta cantidad mide cuánto se mueve el equilibrio instantáneo entre dos tiempos
consecutivos.
\end{definicion}

\begin{definicion}[Coeficiente de contracción de Dobrushin]
Para \(K\in\mathcal S_n\), definimos
\[
\delta(K)
:=
\frac12\max_{1\leq i,r\leq n}
\sum_{j=1}^n |K_{ij}-K_{rj}|.
\]
\end{definicion}

\begin{proposicion}[Descomposición del lag: mezcla contra objetivo móvil]
Supongamos que cada \(K_t\) admite estacionaria instantánea única
\(\pi_t\). Entonces
\[
\ell_{t+1}
\leq
\delta(K_t)\ell_t+
\omega_t.
\]
\end{proposicion}

\begin{proof}
Como \(\pi_tK_t=\pi_t\), se tiene
\begin{align*}
\ell_{t+1}
&=\|X_{t+1}-\pi_{t+1}\|_1\\
&=\|X_tK_t-\pi_{t+1}\|_1\\
&\leq \|X_tK_t-\pi_tK_t\|_1+\|\pi_t-\pi_{t+1}\|_1.
\end{align*}
Por la propiedad contractiva del coeficiente de Dobrushin,
\[
\|X_tK_t-\pi_tK_t\|_1
\leq
\delta(K_t)\|X_t-\pi_t\|_1.
\]
Sustituyendo las definiciones de \(\ell_t\) y \(\omega_t\), obtenemos
\[
\ell_{t+1}
\leq
\delta(K_t)\ell_t+
\omega_t.
\]
\end{proof}

\begin{observacion}[Lectura del segundo experimento]
El experimento extremo fuerza una situación donde el mecanismo \(K_t\) cambia
rápidamente y el equilibrio instantáneo \(\pi_t\) se mueve. Por eso el estado
\(X_t\) puede parecer que ``persigue'' a \(\pi_t\), pero no coincide
necesariamente con ella. La desigualdad
\[
\ell_{t+1}\leq \delta(K_t)\ell_t+\omega_t
\]
separa dos efectos: la capacidad de mezcla de \(K_t\), medida por
\(\delta(K_t)\), y el movimiento del objetivo instantáneo, medido por
\(\omega_t\).
\end{observacion}

\begin{proposicion}[El objetivo móvil puede impedir convergencia]
Existe una dinámica markoviana no autónoma con matrices estocásticas primitivas
para la cual \(X_t\) no converge.
\end{proposicion}

\begin{proof}
Tomemos \(n=2\) y dos distribuciones distintas
\[
a=(0.9,0.1),
\qquad
b=(0.1,0.9).
\]
Definamos matrices de rango uno
\[
K^{a}:=
\begin{pmatrix}
0.9 & 0.1\\
0.9 & 0.1
\end{pmatrix},
\qquad
K^{b}:=
\begin{pmatrix}
0.1 & 0.9\\
0.1 & 0.9
\end{pmatrix}.
\]
Ambas son estocásticas y positivas, luego primitivas. Sea
\[
K_t=
\begin{cases}
K^{a}, & t\text{ par},\\
K^{b}, & t\text{ impar}.
\end{cases}
\]
Para cualquier \(X\in\Delta_2\), se cumple
\[
XK^{a}=a,
\qquad
XK^{b}=b.
\]
Por tanto, la trayectoria alterna entre \(a\) y \(b\). Como \(a\neq b\), la
sucesión \((X_t)\) no converge. Esto muestra que la existencia de estacionarias
instantáneas no implica convergencia cuando el mecanismo cambia con el tiempo.
\end{proof}

\begin{definicion}[Conmutador temporal]
Para \(t\geq 1\), definimos
\[
C_t:=K_tK_{t-1}-K_{t-1}K_t.
\]
La magnitud asociada es
\[
\chi_t:=\|C_t\|_F.
\]
\end{definicion}

\begin{proposicion}[Interpretación del conmutador temporal]
Para todo estado \(X\in\Delta_n\),
\[
XK_tK_{t-1}-XK_{t-1}K_t
=
XC_t.
\]
En particular, si \(C_t\neq 0\), existe un estado interior
\(X_\varepsilon\in\Delta_n^\circ\) tal que
\[
X_\varepsilon K_tK_{t-1}
\neq
X_\varepsilon K_{t-1}K_t.
\]
\end{proposicion}

\begin{proof}
Por asociatividad del producto matricial,
\[
XK_tK_{t-1}-XK_{t-1}K_t
=
X(K_tK_{t-1}-K_{t-1}K_t)
=
XC_t.
\]
Si \(C_t\neq 0\), alguna fila de \(C_t\) es no nula. Sea \(e_i\) un vértice
canónico tal que \(e_iC_t\neq 0\). Definimos
\[
U=\left(\frac1n,\dots,\frac1n\right),
\qquad
X_\varepsilon=(1-\varepsilon)e_i+\varepsilon U.
\]
Para \(\varepsilon\in(0,1)\), se tiene \(X_\varepsilon\in\Delta_n^\circ\). Como
\[
X_\varepsilon C_t
=
(1-\varepsilon)e_iC_t+
\varepsilon UC_t,
\]
y \(e_iC_t\neq 0\), por continuidad existe \(\varepsilon>0\) suficientemente
pequeño tal que
\[
X_\varepsilon C_t\neq 0.
\]
Por tanto,
\[
X_\varepsilon K_tK_{t-1}
\neq
X_\varepsilon K_{t-1}K_t.
\]
\end{proof}

\begin{proposicion}[Cota entre velocidad matricial y no conmutatividad]
Para \(t\geq 1\),
\[
\chi_t
\leq
2\|K_{t-1}\|_2\,\nu_t.
\]
En particular, como \(K_{t-1}\in\mathcal S_n\), se tiene la cota general
\[
\chi_t
\leq
2\sqrt n\,\nu_t.
\]
\end{proposicion}

\begin{proof}
Escribimos
\[
K_tK_{t-1}-K_{t-1}K_t
=
(K_t-K_{t-1})K_{t-1}-K_{t-1}(K_t-K_{t-1}).
\]
Entonces, usando submultiplicatividad de normas,
\begin{align*}
\chi_t
&\leq
\|(K_t-K_{t-1})K_{t-1}\|_F
+
\|K_{t-1}(K_t-K_{t-1})\|_F\\
&\leq
\|K_t-K_{t-1}\|_F\|K_{t-1}\|_2
+
\|K_{t-1}\|_2\|K_t-K_{t-1}\|_F\\
&=
2\|K_{t-1}\|_2\nu_t.
\end{align*}
Finalmente, para una matriz estocástica por filas, cada fila tiene norma
\(2\) menor o igual que \(1\), luego
\[
\|K_{t-1}\|_F\leq \sqrt n.
\]
Como \(\|K_{t-1}\|_2\leq\|K_{t-1}\|_F\), se obtiene
\[
\chi_t\leq 2\sqrt n\,\nu_t.
\]
\end{proof}

\begin{observacion}[Alcance formal de los dos scripts]
Los dos scripts no demuestran todavía una teoría general de control informacional.
Lo que sí verifican formalmente es lo siguiente:
\begin{enumerate}[label=\textup{(\roman*)}]
\item es posible hacer evolucionar la matriz estocástica \(K_t\) sin salir de
\(\mathcal S_n\);
\item el estado \(X_t\) permanece en el símplex mientras \(K_t\) sea estocástica;
\item la relajación convexa hacia \(K_\ast\) tiene convergencia geométrica
explícita;
\item en una dinámica extrema proyectada, el estado puede perseguir una
estacionaria móvil sin converger;
\item la no conmutatividad temporal mide una diferencia real entre aplicar
mecanismos en órdenes distintos.
\end{enumerate}
Por tanto, el objeto matemático pertinente en esta etapa no es una única matriz
de Markov fija, sino una trayectoria de mecanismos
\[
K_0,K_1,K_2,\dots
\]
acoplada a la trayectoria de estados
\[
X_0,X_1,X_2,\dots.
\]
\end{observacion}

% ============================================================
% CONJETURAS PARA PROYECTO II
% (no forman parte del texto; son líneas de investigación abiertas)
% ============================================================

% ------------------------------------------------------------
% CONJETURA A (Ruta 1) — El conmutador acumulado controla
%   la diferencia entre órdenes de aplicación.
%
% Sean K_0,...,K_{t-1} \in S_n y defínase
%   K_{[0,t-1]}^{->}  = K_0 K_1 ... K_{t-1}  (orden cronológico)
%   K_{[0,t-1]}^{<-}  = K_{t-1} ... K_1 K_0  (orden inverso)
% Conjetura: existe una constante C(n) tal que para todo
%   X_0 \in \Delta_n,
%
%   \| X_0 K_{[0,t-1]}^{->} - X_0 K_{[0,t-1]}^{<-} \|_1
%       \leq C(n) \sum_{s=1}^{t-1} \chi_s.
%
% En particular, si \chi_s = 0 para todo s (todas las matrices
% consecutivas conmutan), los dos productos coinciden para
% cualquier estado inicial.
% ------------------------------------------------------------

% ------------------------------------------------------------
% CONJETURA B (Ruta 2) — La presión de seguimiento \omega_t
%   es controlable por la velocidad matricial \nu_t.
%
% Si K_t es primitiva con brecha espectral 1 - \lambda_2(K_t),
% entonces existe C(K_t) > 0 tal que
%
%   \omega_t := \| \pi_{t+1} - \pi_t \|_1
%             \leq \frac{\nu_t}{1 - \lambda_2(K_{t-1})}.
%
% Consecuencia: la descomposición del lag se vuelve un sistema
% cerrado en términos de observables controlables,
%
%   \ell_{t+1} \leq \delta(K_t)\,\ell_t
%               + \frac{\nu_t}{1 - \lambda_2(K_{t-1})},
%
% lo que permite formular el problema de control óptimo:
%
%   \min_{(K_t) \subset S_n}
%       \sum_{t=0}^{T-1} \bigl( \ell_t^2 + \lambda\,\nu_t^2 \bigr)
%   \quad\text{s.a.}\quad
%   \ell_{t+1} \leq \delta(K_t)\,\ell_t + \omega_t.
% ------------------------------------------------------------

% ------------------------------------------------------------
% CONJETURA C (Ruta 3) — La composición Bayes-Markov no es
%   ni puramente local ni puramente global.
%
% Sea B_\mu(X) la actualización bayesiana con verosimilitud \mu
% (operador global del grupo G) y sea L_P el operador de Markov
% (familia de correcciones locales).  Conjetura:
%
%   L_P \circ B_\mu \neq B_\nu \circ L_Q
%
% para ningún par (\nu, Q) en general, es decir, la composición
% no puede reordenarse libremente.  El obstáculo exacto es la
% sección X \mapsto Y^P(X): la corrección local de Markov
% evaluada en B_\mu(X) depende de X a través de \mu, mientras
% que la corrección global de Bayes no depende de X.
%
% Un resultado más fino: caracterizar el conjunto de pares
% (\mu, P) para los cuales sí existe (\nu, Q) con
% L_P \circ B_\mu = B_\nu \circ L_Q.  La condición esperada
% involucra que P tenga filas proporcionales a \mu.
% ------------------------------------------------------------


\clearpage
\section{Bayes y control informacional}

\begin{definicion}[Evidencia bayesiana positiva]
Sea \(n\geq 2\). Una \emph{evidencia bayesiana positiva} sobre \(n\)
hipótesis es un vector
\[
L=(L_1,\dots,L_n)\in \mathbb{R}_{>0}^{n}.
\]

Cada coordenada \(L_i\) representa el peso de verosimilitud asociado a la
hipótesis \(i\). La condición
\[
L_i>0
\qquad
\text{para todo } i=1,\dots,n
\]
garantiza que ninguna hipótesis con masa previa positiva sea anulada por la
actualización.
\end{definicion}

\begin{definicion}[Operador bayesiano]
Sea \(X\in\Delta_n^\circ\) una distribución a priori y sea
\(L\in\mathbb{R}_{>0}^{n}\) una evidencia bayesiana positiva. Se define el
\emph{operador bayesiano asociado a \(L\)} como la aplicación
\[
B_L:\Delta_n^\circ\longrightarrow \Delta_n^\circ
\]
dada por
\[
B_L(X)_i
:=
\frac{X_iL_i}{\sum_{k=1}^{n}X_kL_k},
\qquad
i=1,\dots,n.
\]

Al escalar
\[
Z_L(X)
:=
\sum_{k=1}^{n}X_kL_k
=
\langle X,L\rangle
\]
se le llamará \emph{constante de normalización bayesiana}.
\end{definicion}



\begin{proposicion}[Preservación del interior bajo actualización bayesiana]
Sea \(L\in\mathbb{R}_{>0}^{n}\) una evidencia bayesiana positiva. Entonces,
para todo \(X\in\Delta_n^\circ\), se cumple
\[
B_L(X)\in\Delta_n^\circ.
\]
En particular, el operador bayesiano está bien definido como aplicación
\[
B_L:\Delta_n^\circ\longrightarrow \Delta_n^\circ.
\]
\end{proposicion}

\begin{proof}
Sea \(X\in\Delta_n^\circ\). Entonces
\[
X_i>0
\qquad
\text{para todo } i=1,\dots,n.
\]
Como \(L\in\mathbb{R}_{>0}^{n}\), también se tiene
\[
L_i>0
\qquad
\text{para todo } i=1,\dots,n.
\]
Por tanto,
\[
X_iL_i>0
\qquad
\text{para todo } i=1,\dots,n.
\]

Además, la constante de normalización satisface
\[
Z_L(X)
=
\sum_{k=1}^{n}X_kL_k.
\]
Como cada sumando \(X_kL_k\) es estrictamente positivo, se obtiene
\[
Z_L(X)>0.
\]
Así, cada coordenada de \(B_L(X)\) está bien definida y cumple
\[
B_L(X)_i
=
\frac{X_iL_i}{Z_L(X)}
>
0.
\]

Resta verificar la normalización. En efecto,
\[
\sum_{i=1}^{n}B_L(X)_i
=
\sum_{i=1}^{n}
\frac{X_iL_i}{Z_L(X)}
=
\frac{1}{Z_L(X)}
\sum_{i=1}^{n}X_iL_i.
\]
Por la definición de \(Z_L(X)\),
\[
\sum_{i=1}^{n}B_L(X)_i
=
\frac{Z_L(X)}{Z_L(X)}
=
1.
\]

Por tanto, \(B_L(X)\) tiene coordenadas estrictamente positivas y suma \(1\). Luego
\[
B_L(X)\in\Delta_n^\circ.
\]
\end{proof}

\begin{observacion}[Razón de la positividad estricta]
La condición
\[
L\in\mathbb{R}_{>0}^{n}
\]
no es solamente técnica. Si alguna coordenada \(L_i\) fuera cero, entonces aun cuando
\(X_i>0\), se tendría
\[
B_L(X)_i=0,
\]
y el posterior podría salir del interior del símplex.

Por eso, dentro de esta sección se trabaja con verosimilitudes estrictamente positivas:
la actualización bayesiana conserva el soporte completo del estado inicial y permanece en
\(\Delta_n^\circ\).
\end{observacion}

\begin{proposicion}[Invariancia por escala de la evidencia]
Sea \(L\in\mathbb{R}_{>0}^{n}\) una evidencia bayesiana positiva y sea
\(c>0\). Entonces
\[
B_{cL}=B_L.
\]
Es decir, para todo \(X\in\Delta_n^\circ\),
\[
B_{cL}(X)=B_L(X).
\]
\end{proposicion}

\begin{proof}
Sea \(X\in\Delta_n^\circ\). Para cada \(i=1,\dots,n\), se tiene
\[
B_{cL}(X)_i
=
\frac{X_i(cL_i)}
{\sum_{k=1}^{n}X_k(cL_k)}.
\]
Factorizando \(c\) en el numerador y en el denominador,
\[
B_{cL}(X)_i
=
\frac{cX_iL_i}
{c\sum_{k=1}^{n}X_kL_k}.
\]
Como \(c>0\), se cancela el factor común:
\[
B_{cL}(X)_i
=
\frac{X_iL_i}
{\sum_{k=1}^{n}X_kL_k}
=
B_L(X)_i.
\]
Por tanto,
\[
B_{cL}(X)_i=B_L(X)_i
\qquad
\text{para todo } i=1,\dots,n.
\]
Luego
\[
B_{cL}(X)=B_L(X).
\]
Como \(X\in\Delta_n^\circ\) fue arbitrario, se concluye que
\[
B_{cL}=B_L.
\]
\end{proof}

\begin{observacion}[La evidencia es proyectiva]
La proposición anterior muestra que la evidencia bayesiana no determina un operador por su
magnitud absoluta, sino por sus proporciones internas.

En particular,
\[
L
\quad\text{y}\quad
cL,
\qquad c>0,
\]
representan la misma actualización bayesiana sobre \(\Delta_n^\circ\).

Por tanto, el objeto estructural no es estrictamente el vector \(L\in\mathbb{R}_{>0}^{n}\),
sino su clase bajo reescalamientos positivos globales.
\end{observacion}


\begin{definicion}[Clase proyectiva de una evidencia]
En \(\mathbb{R}_{>0}^{n}\) se define la relación
\[
L\sim M
\]
si y sólo si existe una constante \(c>0\) tal que
\[
M=cL.
\]

La clase de equivalencia de \(L\) será denotada por
\[
[L]
:=
\{cL:c>0\}.
\]

Al cociente
\[
\mathbb{P}^{n-1}_{+}
:=
\mathbb{R}_{>0}^{n}/\sim
\]
lo llamaremos \emph{espacio proyectivo positivo de evidencias}.
\end{definicion}

\begin{proposicion}[Identificación proyectiva de operadores bayesianos]
Sean \(L,M\in\mathbb{R}_{>0}^{n}\). Entonces
\[
B_L=B_M
\]
si y sólo si
\[
[L]=[M].
\]
Equivalente, dos evidencias positivas inducen el mismo operador bayesiano si y sólo si
difieren por un factor escalar positivo global.
\end{proposicion}

\begin{proof}
Supongamos primero que
\[
[L]=[M].
\]
Entonces existe \(c>0\) tal que
\[
M=cL.
\]
Por la invariancia por escala de la evidencia,
\[
B_M=B_{cL}=B_L.
\]

Recíprocamente, supongamos que
\[
B_L=B_M.
\]
Sea \(X\in\Delta_n^\circ\). Entonces, para cada \(i=1,\dots,n\),
\[
B_L(X)_i=B_M(X)_i.
\]
Por definición,
\[
\frac{X_iL_i}{\sum_{k=1}^{n}X_kL_k}
=
\frac{X_iM_i}{\sum_{k=1}^{n}X_kM_k}.
\]
Como \(X_i>0\), podemos cancelar \(X_i\) y obtener
\[
\frac{L_i}{\sum_{k=1}^{n}X_kL_k}
=
\frac{M_i}{\sum_{k=1}^{n}X_kM_k}.
\]
Por tanto,
\[
M_i
=
\frac{\sum_{k=1}^{n}X_kM_k}
{\sum_{k=1}^{n}X_kL_k}
L_i.
\]
Definiendo
\[
c
:=
\frac{\sum_{k=1}^{n}X_kM_k}
{\sum_{k=1}^{n}X_kL_k},
\]
se tiene \(c>0\) y
\[
M_i=cL_i
\qquad
\text{para todo } i=1,\dots,n.
\]
Así,
\[
M=cL,
\]
y por tanto
\[
[L]=[M].
\]
\end{proof}



\begin{definicion}[Producto de evidencias]
Sean \(L^{(1)},L^{(2)}\in\mathbb{R}_{>0}^{n}\) dos evidencias bayesianas
positivas. Definimos su producto coordenado por
\[
L^{(1)}\odot L^{(2)}
:=
\bigl(
L^{(1)}_1L^{(2)}_1,\dots,
L^{(1)}_nL^{(2)}_n
\bigr).
\]
Como cada coordenada es estrictamente positiva, se tiene
\[
L^{(1)}\odot L^{(2)}\in\mathbb{R}_{>0}^{n}.
\]
\end{definicion}

\begin{proposicion}[Composición de actualizaciones bayesianas]
Sean \(L^{(1)},L^{(2)}\in\mathbb{R}_{>0}^{n}\). Entonces
\[
B_{L^{(2)}}\circ B_{L^{(1)}}
=
B_{L^{(1)}\odot L^{(2)}}.
\]
En particular, la composición de dos operadores bayesianos vuelve a ser un
operador bayesiano.
\end{proposicion}

\begin{proof}
Sea \(X\in\Delta_n^\circ\) y escribamos
\[
X' := B_{L^{(1)}}(X).
\]
Entonces
\[
X'_i
=
\frac{X_iL^{(1)}_i}
{\sum_{j=1}^{n}X_jL^{(1)}_j}.
\]
Ahora aplicamos \(B_{L^{(2)}}\) a \(X'\):
\[
B_{L^{(2)}}(X')_i
=
\frac{X'_iL^{(2)}_i}
{\sum_{k=1}^{n}X'_kL^{(2)}_k}.
\]
Sustituyendo la expresión de \(X'\), obtenemos
\[
B_{L^{(2)}}(X')_i
=
\frac{
\frac{X_iL^{(1)}_i}{\sum_{j=1}^{n}X_jL^{(1)}_j}
L^{(2)}_i
}{
\sum_{k=1}^{n}
\frac{X_kL^{(1)}_k}{\sum_{j=1}^{n}X_jL^{(1)}_j}
L^{(2)}_k
}.
\]
El factor común
\[
\sum_{j=1}^{n}X_jL^{(1)}_j
\]
se cancela en numerador y denominador. Por tanto,
\[
B_{L^{(2)}}(X')_i
=
\frac{X_iL^{(1)}_iL^{(2)}_i}
{\sum_{k=1}^{n}X_kL^{(1)}_kL^{(2)}_k}.
\]
Por definición del producto coordenado,
\[
L^{(1)}_iL^{(2)}_i
=
\bigl(L^{(1)}\odot L^{(2)}\bigr)_i.
\]
Así,
\[
B_{L^{(2)}}(B_{L^{(1)}}(X))_i
=
B_{L^{(1)}\odot L^{(2)}}(X)_i.
\]
Como esto vale para toda coordenada \(i\), se concluye que
\[
B_{L^{(2)}}\circ B_{L^{(1)}}
=
B_{L^{(1)}\odot L^{(2)}}.
\]
\end{proof}


\begin{definicion}[Trayectoria bayesiana]
Sea \(X_0\in\Delta_n^\circ\) un estado inicial y sea
\[
\{L^{(t)}\}_{t\geq 1}\subset \mathbb{R}_{>0}^{n}
\]
una sucesión de evidencias bayesianas positivas.

La sucesión \(\{X_t\}_{t\geq 0}\subset\Delta_n^\circ\) definida por
\[
X_{t+1}
=
B_{L^{(t+1)}}(X_t),
\qquad t\geq 0,
\]
se llamará \emph{trayectoria bayesiana generada por}
\[
X_0,\; L^{(1)},L^{(2)},\dots.
\]
\end{definicion}

\begin{proposicion}[Forma cerrada de una trayectoria bayesiana]
Sea \(\{X_t\}_{t\geq 0}\) una trayectoria bayesiana generada por
\(X_0\in\Delta_n^\circ\) y por evidencias
\[
L^{(1)},L^{(2)},\dots,L^{(t)}
\in\mathbb{R}_{>0}^{n}.
\]
Defínase, para cada \(i=1,\dots,n\),
\[
\Lambda_i^{(t)}
:=
\prod_{s=1}^{t}L_i^{(s)}.
\]
Entonces, para todo \(t\geq 1\),
\[
X_{t,i}
=
\frac{
X_{0,i}\Lambda_i^{(t)}
}{
\sum_{k=1}^{n}X_{0,k}\Lambda_k^{(t)}
},
\qquad i=1,\dots,n.
\]
\end{proposicion}

\begin{proof}
Procedemos por inducción sobre \(t\).

Para \(t=1\), por definición de la actualización bayesiana,
\[
X_{1,i}
=
B_{L^{(1)}}(X_0)_i
=
\frac{X_{0,i}L_i^{(1)}}{\sum_{k=1}^{n}X_{0,k}L_k^{(1)}}.
\]
Como
\[
\Lambda_i^{(1)}=L_i^{(1)},
\]
la fórmula se cumple para \(t=1\).

Supongamos ahora que la fórmula es válida para algún \(t\geq 1\), es decir,
\[
X_{t,i}
=
\frac{
X_{0,i}\Lambda_i^{(t)}
}{
\sum_{k=1}^{n}X_{0,k}\Lambda_k^{(t)}
}.
\]
Entonces
\[
X_{t+1,i}
=
B_{L^{(t+1)}}(X_t)_i
=
\frac{X_{t,i}L_i^{(t+1)}}{\sum_{m=1}^{n}X_{t,m}L_m^{(t+1)}}.
\]
Sustituyendo la hipótesis inductiva,
\[
X_{t+1,i}
=
\frac{
\frac{X_{0,i}\Lambda_i^{(t)}}{\sum_{k=1}^{n}X_{0,k}\Lambda_k^{(t)}}
L_i^{(t+1)}
}{
\sum_{m=1}^{n}
\frac{X_{0,m}\Lambda_m^{(t)}}{\sum_{k=1}^{n}X_{0,k}\Lambda_k^{(t)}}
L_m^{(t+1)}
}.
\]
El factor común
\[
\sum_{k=1}^{n}X_{0,k}\Lambda_k^{(t)}
\]
se cancela. Por tanto,
\[
X_{t+1,i}
=
\frac{
X_{0,i}\Lambda_i^{(t)}L_i^{(t+1)}
}{
\sum_{m=1}^{n}
X_{0,m}\Lambda_m^{(t)}L_m^{(t+1)}
}.
\]
Como
\[
\Lambda_i^{(t+1)}
=
\Lambda_i^{(t)}L_i^{(t+1)},
\]
obtenemos
\[
X_{t+1,i}
=
\frac{
X_{0,i}\Lambda_i^{(t+1)}
}{
\sum_{m=1}^{n}X_{0,m}\Lambda_m^{(t+1)}
}.
\]
Por inducción, la fórmula vale para todo \(t\geq 1\).
\end{proof}


\begin{definicion}[Log-verosimilitud acumulada]
Sea
\[
\{L^{(s)}\}_{s=1}^{t}\subset \mathbb{R}_{>0}^{n}
\]
una familia finita de evidencias bayesianas positivas. Se define la
\emph{log-verosimilitud acumulada} hasta el tiempo \(t\) como el vector
\[
\theta^{(t)}=(\theta_1^{(t)},\dots,\theta_n^{(t)})\in\mathbb{R}^{n}
\]
dado por
\[
\theta_i^{(t)}
:=
\sum_{s=1}^{t}\log L_i^{(s)},
\qquad
i=1,\dots,n.
\]

Equivalentemente,
\[
\theta_i^{(t)}
=
\log\left(\prod_{s=1}^{t}L_i^{(s)}\right)
=
\log \Lambda_i^{(t)}.
\]
\end{definicion}

\begin{proposicion}[Trayectoria bayesiana en coordenadas logarítmicas]
Sea \(\{X_t\}_{t\geq 0}\) una trayectoria bayesiana generada por
\(X_0\in\Delta_n^\circ\) y evidencias positivas
\[
L^{(1)},L^{(2)},\dots,L^{(t)}\in\mathbb{R}_{>0}^{n}.
\]
Si
\[
\theta_i^{(t)}
=
\sum_{s=1}^{t}\log L_i^{(s)},
\]
entonces
\[
X_{t,i}
=
\frac{
X_{0,i}\exp(\theta_i^{(t)})
}{
\sum_{k=1}^{n}X_{0,k}\exp(\theta_k^{(t)})
},
\qquad
i=1,\dots,n.
\]
En particular, la composición de actualizaciones bayesianas corresponde a la suma de
log-verosimilitudes.
\end{proposicion}

\begin{proof}
Por la forma cerrada de la trayectoria bayesiana, se tiene
\[
X_{t,i}
=
\frac{
X_{0,i}\Lambda_i^{(t)}
}{
\sum_{k=1}^{n}X_{0,k}\Lambda_k^{(t)}
},
\]
donde
\[
\Lambda_i^{(t)}
=
\prod_{s=1}^{t}L_i^{(s)}.
\]
Por definición de \(\theta_i^{(t)}\),
\[
\theta_i^{(t)}
=
\log \Lambda_i^{(t)}.
\]
Como \(\Lambda_i^{(t)}>0\), se sigue que
\[
\Lambda_i^{(t)}
=
\exp(\theta_i^{(t)}).
\]
Sustituyendo en la fórmula cerrada,
\[
X_{t,i}
=
\frac{
X_{0,i}\exp(\theta_i^{(t)})
}{
\sum_{k=1}^{n}X_{0,k}\exp(\theta_k^{(t)})
}.
\]

Finalmente, si se agregan dos bloques de evidencia acumulada,
\[
\theta^{(a)}
=
\sum_{s=1}^{a}\log L^{(s)}
\]
y
\[
\theta^{(b)}
=
\sum_{s=a+1}^{a+b}\log L^{(s)},
\]
entonces la evidencia total satisface
\[
\theta^{(a+b)}
=
\theta^{(a)}+\theta^{(b)}.
\]
Por tanto, la composición multiplicativa de evidencias se transforma en suma en
coordenadas logarítmicas.
\end{proof}

\begin{observacion}[Lectura aditiva de la información]
En coordenadas de evidencia directa, la información se acumula mediante productos:
\[
\Lambda^{(t)}
=
L^{(1)}\odot\cdots\odot L^{(t)}.
\]
En coordenadas logarítmicas, la misma acumulación se escribe aditivamente:
\[
\theta^{(t)}
=
\sum_{s=1}^{t}\log L^{(s)}.
\]

Esta es la forma más conveniente para conectar la actualización bayesiana con la ley de
composición de los operadores globales de Control Informacional.
\end{observacion}



\begin{definicion}[Evidencia dirigida hacia un objetivo]
Sean \(X,\pi\in\Delta_n^\circ\). Se define la \emph{evidencia dirigida de
\(X\) hacia \(\pi\)} como el vector
\[
L^{X\to \pi}
\in \mathbb{R}_{>0}^{n}
\]
dado por
\[
L_i^{X\to \pi}
:=
\frac{\pi_i}{X_i},
\qquad
i=1,\dots,n.
\]

Equivalentemente, se define la corrección global asociada
\[
Y^{X\to \pi}\in C_{\mathrm{glob}}
\]
por
\[
1+Y_i^{X\to \pi}
:=
\frac{\pi_i}{X_i},
\]
es decir,
\[
Y_i^{X\to \pi}
=
\frac{\pi_i}{X_i}-1.
\]
\end{definicion}

\begin{proposicion}[Actualización exacta hacia un objetivo]
Sean \(X,\pi\in\Delta_n^\circ\). Entonces la evidencia dirigida
\(L^{X\to \pi}\) satisface
\[
B_{L^{X\to \pi}}(X)=\pi.
\]
Equivalentemente,
\[
T_{Y^{X\to \pi}}(X)=\pi.
\]
\end{proposicion}

\begin{proof}
Como \(X,\pi\in\Delta_n^\circ\), se tiene
\[
X_i>0
\qquad\text{y}\qquad
\pi_i>0
\]
para todo \(i=1,\dots,n\). Por tanto,
\[
L_i^{X\to \pi}
=
\frac{\pi_i}{X_i}
>
0,
\]
y así
\[
L^{X\to \pi}\in\mathbb{R}_{>0}^{n}.
\]

Ahora calculamos la actualización bayesiana inducida por esta evidencia. Para cada
\(i=1,\dots,n\),
\[
B_{L^{X\to \pi}}(X)_i
=
\frac{
X_iL_i^{X\to \pi}
}{
\sum_{k=1}^{n}X_kL_k^{X\to \pi}
}.
\]
Sustituyendo
\[
L_i^{X\to \pi}
=
\frac{\pi_i}{X_i},
\]
obtenemos
\[
B_{L^{X\to \pi}}(X)_i
=
\frac{
X_i\frac{\pi_i}{X_i}
}{
\sum_{k=1}^{n}
X_k\frac{\pi_k}{X_k}
}.
\]
Por cancelación,
\[
B_{L^{X\to \pi}}(X)_i
=
\frac{
\pi_i
}{
\sum_{k=1}^{n}\pi_k
}.
\]
Como \(\pi\in\Delta_n^\circ\), se tiene
\[
\sum_{k=1}^{n}\pi_k=1.
\]
Luego
\[
B_{L^{X\to \pi}}(X)_i=\pi_i.
\]
Esto vale para toda coordenada \(i\), por lo que
\[
B_{L^{X\to \pi}}(X)=\pi.
\]

Finalmente, como
\[
1+Y_i^{X\to \pi}=L_i^{X\to \pi},
\]
por la identificación entre operadores bayesianos y operadores globales de Control
Informacional se obtiene también
\[
T_{Y^{X\to \pi}}(X)=\pi.
\]
\end{proof}

\begin{proposicion}[Unicidad proyectiva de la evidencia dirigida]
Sean \(X,\pi\in\Delta_n^\circ\). Si \(M\in\mathbb{R}_{>0}^{n}\) satisface
\[
B_M(X)=\pi,
\]
entonces
\[
[M]=[L^{X\to \pi}].
\]
Equivalente, existe una constante \(c>0\) tal que
\[
M_i
=
c\,\frac{\pi_i}{X_i},
\qquad
i=1,\dots,n.
\]
\end{proposicion}

\begin{proof}
Supongamos que
\[
B_M(X)=\pi.
\]
Entonces, para cada \(i=1,\dots,n\),
\[
\frac{X_iM_i}{\sum_{k=1}^{n}X_kM_k}
=
\pi_i.
\]
Definamos
\[
Z_M(X)
:=
\sum_{k=1}^{n}X_kM_k.
\]
Como \(X\in\Delta_n^\circ\) y \(M\in\mathbb{R}_{>0}^{n}\), se tiene
\[
Z_M(X)>0.
\]
De la igualdad anterior se obtiene
\[
X_iM_i
=
\pi_i Z_M(X).
\]
Como \(X_i>0\), podemos despejar:
\[
M_i
=
Z_M(X)\frac{\pi_i}{X_i}.
\]
Tomando
\[
c:=Z_M(X)>0,
\]
queda
\[
M_i
=
c\,L_i^{X\to \pi}.
\]
Por tanto,
\[
M=cL^{X\to \pi},
\]
y en consecuencia
\[
[M]=[L^{X\to \pi}].
\]
\end{proof}

\begin{observacion}[Normalización canónica]
La evidencia dirigida
\[
L_i^{X\to \pi}
=
\frac{\pi_i}{X_i}
\]
es un representante canónico de la clase proyectiva de evidencias que llevan \(X\) hacia
\(\pi\). En efecto,
\[
\langle X,L^{X\to \pi}\rangle
=
\sum_{i=1}^{n}X_i\frac{\pi_i}{X_i}
=
\sum_{i=1}^{n}\pi_i
=
1.
\]

Por tanto, aunque todo vector de la forma
\[
cL^{X\to \pi},
\qquad c>0,
\]
produce la misma actualización, el representante \(L^{X\to \pi}\) queda distinguido por la
normalización
\[
\langle X,L^{X\to \pi}\rangle=1.
\]
\end{observacion}

\begin{definicion}[Campo infinitesimal bayesiano]
Sea \(v\in\mathbb{R}^{n}\). Para cada \(\varepsilon\in\mathbb{R}\), definimos la evidencia
positiva
\[
L^{(\varepsilon,v)}
:=
\exp(\varepsilon v)
=
\bigl(
e^{\varepsilon v_1},\dots,e^{\varepsilon v_n}
\bigr).
\]
El flujo bayesiano asociado a \(v\) se define por
\[
\Phi_\varepsilon^v(X)
:=
B_{L^{(\varepsilon,v)}}(X),
\qquad X\in\Delta_n^\circ.
\]
El \emph{campo infinitesimal bayesiano} generado por \(v\) es
\[
V_B^v(X)
:=
\left.\frac{d}{d\varepsilon}\right|_{\varepsilon=0}
\Phi_\varepsilon^v(X).
\]
\end{definicion}

\begin{proposicion}[Forma explícita del campo infinitesimal bayesiano]
Sea \(v\in\mathbb{R}^{n}\). Entonces, para todo \(X\in\Delta_n^\circ\),
\[
\bigl(V_B^v(X)\bigr)_i
=
X_i\left(v_i-\sum_{k=1}^{n}X_kv_k\right),
\qquad i=1,\dots,n.
\]
En particular,
\[
V_B^v(X)\in T_X\Delta_n^\circ.
\]
\end{proposicion}

\begin{proof}
Por definición,
\[
\Phi_\varepsilon^v(X)_i
=
B_{\exp(\varepsilon v)}(X)_i
=
\frac{
X_i e^{\varepsilon v_i}
}{
\sum_{k=1}^{n}X_k e^{\varepsilon v_k}
}.
\]
Escribamos
\[
Z(\varepsilon)
:=
\sum_{k=1}^{n}X_k e^{\varepsilon v_k}.
\]
Entonces
\[
\Phi_\varepsilon^v(X)_i
=
\frac{X_i e^{\varepsilon v_i}}{Z(\varepsilon)}.
\]
Derivando respecto a \(\varepsilon\),
\[
\frac{d}{d\varepsilon}\Phi_\varepsilon^v(X)_i
=
\frac{
X_i v_i e^{\varepsilon v_i}Z(\varepsilon)
-
X_i e^{\varepsilon v_i}Z'(\varepsilon)
}{
Z(\varepsilon)^2
}.
\]
Ahora,
\[
Z(0)=\sum_{k=1}^{n}X_k=1
\]
y
\[
Z'(0)
=
\sum_{k=1}^{n}X_kv_k.
\]
Evaluando en \(\varepsilon=0\), obtenemos
\[
\left.\frac{d}{d\varepsilon}\right|_{\varepsilon=0}
\Phi_\varepsilon^v(X)_i
=
X_i
\left(
v_i-\sum_{k=1}^{n}X_kv_k
\right).
\]
Por tanto,
\[
\bigl(V_B^v(X)\bigr)_i
=
X_i
\left(
v_i-\sum_{k=1}^{n}X_kv_k
\right).
\]

Finalmente,
\[
\sum_{i=1}^{n}\bigl(V_B^v(X)\bigr)_i
=
\sum_{i=1}^{n}X_iv_i
-
\left(\sum_{k=1}^{n}X_kv_k\right)
\sum_{i=1}^{n}X_i.
\]
Como
\[
\sum_{i=1}^{n}X_i=1,
\]
se sigue que
\[
\sum_{i=1}^{n}\bigl(V_B^v(X)\bigr)_i
=
\sum_{i=1}^{n}X_iv_i
-
\sum_{k=1}^{n}X_kv_k
=
0.
\]
Luego
\[
V_B^v(X)\in T_X\Delta_n^\circ.
\]
\end{proof}

\begin{proposicion}[Interpretación Fisher del campo infinitesimal bayesiano]
Sea \(v\in\mathbb{R}^{n}\) y defínase
\[
F_v:\Delta_n^\circ\longrightarrow\mathbb{R},
\qquad
F_v(X):=\langle X,v\rangle
=
\sum_{i=1}^{n}X_iv_i.
\]
Entonces
\[
V_B^v(X)=\nabla_g F_v(X),
\]
donde \(\nabla_g\) denota el gradiente respecto a la métrica de Fisher
\[
g_X(u,w)
=
\sum_{i=1}^{n}\frac{u_iw_i}{X_i}.
\]
\end{proposicion}

\begin{proof}
Sea \(u\in T_X\Delta_n^\circ\). Entonces
\[
\sum_{i=1}^{n}u_i=0.
\]
Por la fórmula explícita del campo infinitesimal bayesiano,
\[
\bigl(V_B^v(X)\bigr)_i
=
X_i
\left(
v_i-\sum_{k=1}^{n}X_kv_k
\right).
\]
Luego
\[
g_X(V_B^v(X),u)
=
\sum_{i=1}^{n}
\frac{
X_i\left(v_i-\sum_{k=1}^{n}X_kv_k\right)u_i
}{X_i}.
\]
Cancelando \(X_i\), obtenemos
\[
g_X(V_B^v(X),u)
=
\sum_{i=1}^{n}v_iu_i
-
\left(\sum_{k=1}^{n}X_kv_k\right)
\sum_{i=1}^{n}u_i.
\]
Como \(u\in T_X\Delta_n^\circ\),
\[
\sum_{i=1}^{n}u_i=0.
\]
Por tanto,
\[
g_X(V_B^v(X),u)
=
\sum_{i=1}^{n}v_iu_i.
\]
Pero
\[
dF_v(X)[u]
=
\sum_{i=1}^{n}v_iu_i.
\]
Así,
\[
g_X(V_B^v(X),u)=dF_v(X)[u]
\]
para todo \(u\in T_X\Delta_n^\circ\). En consecuencia,
\[
V_B^v(X)=\nabla_g F_v(X).
\]
\end{proof}

\begin{observacion}[Carácter proyectivo del campo infinitesimal]
Si \(c\in\mathbb{R}\), entonces
\[
V_B^{v+c\mathbf{1}}(X)=V_B^v(X).
\]
En efecto,
\[
(v_i+c)-\sum_{k=1}^{n}X_k(v_k+c)
=
v_i+c-\sum_{k=1}^{n}X_kv_k-c\sum_{k=1}^{n}X_k
=
v_i-\sum_{k=1}^{n}X_kv_k.
\]
Por tanto, el campo infinitesimal bayesiano no depende de \(v\) como vector absoluto,
sino de su clase módulo constantes:
\[
v\sim v+c\mathbf{1}.
\]
Esto es la versión infinitesimal de la invariancia proyectiva de la evidencia bayesiana.
\end{observacion}

\begin{definicion}[Conmutador Bayes--Markov]
Sea \(P\in\mathbb{R}^{n\times n}\) una matriz estocástica por filas estrictamente positiva
y sea \(L\in\mathbb{R}_{>0}^{n}\) una evidencia bayesiana positiva.

Definimos el \emph{conmutador Bayes--Markov} asociado a \(P\) y \(L\) como la aplicación
\[
\mathcal{C}_{P,L}:\Delta_n^\circ\longrightarrow T\Delta_n^\circ
\]
dada por
\[
\mathcal{C}_{P,L}(X)
:=
(B_L\circ L_P)(X)
-
(L_P\circ B_L)(X).
\]
En coordenadas,
\[
\mathcal{C}_{P,L}(X)_j
=
\frac{(XP)_jL_j}{\sum_{m=1}^{n}(XP)_mL_m}
-
\frac{\sum_{i=1}^{n}X_iL_iP_{ij}}{\sum_{k=1}^{n}X_kL_k}.
\]
\end{definicion}

\begin{proposicion}[El conmutador Bayes--Markov es tangente al símplex]
Sea \(P\) una matriz estocástica por filas estrictamente positiva y sea
\(L\in\mathbb{R}_{>0}^{n}\). Entonces, para todo \(X\in\Delta_n^\circ\),
\[
\sum_{j=1}^{n}\mathcal{C}_{P,L}(X)_j=0.
\]
Por tanto,
\[
\mathcal{C}_{P,L}(X)\in T_X\Delta_n^\circ.
\]
\end{proposicion}

\begin{proof}
Como \(P\) es estocástica por filas y estrictamente positiva,
\[
L_P(X)=XP\in\Delta_n^\circ.
\]
Luego
\[
B_L(L_P(X))\in\Delta_n^\circ.
\]
También,
\[
B_L(X)\in\Delta_n^\circ
\]
y, por preservación del símplex bajo \(L_P\),
\[
L_P(B_L(X))\in\Delta_n^\circ.
\]
Por tanto, ambos vectores
\[
(B_L\circ L_P)(X)
\qquad\text{y}\qquad
(L_P\circ B_L)(X)
\]
pertenecen a \(\Delta_n^\circ\), de modo que cada uno tiene suma total igual a \(1\).
Así,
\[
\sum_{j=1}^{n}\mathcal{C}_{P,L}(X)_j
=
\sum_{j=1}^{n}(B_L\circ L_P)(X)_j
-
\sum_{j=1}^{n}(L_P\circ B_L)(X)_j
=
1-1
=
0.
\]
Luego
\[
\mathcal{C}_{P,L}(X)\in T_X\Delta_n^\circ.
\]
\end{proof}


\begin{proposicion}[Condición de conmutación Bayes--Markov]
Sea \(P\in\mathbb{R}^{n\times n}\) una matriz estocástica por filas estrictamente
positiva con distribución estacionaria \(\pi\in\Delta_n^\circ\), y sea
\(L\in\mathbb{R}_{>0}^{n}\). Entonces
\[
B_L\circ L_P = L_P\circ B_L
\]
si y sólo si existe \(c>0\) tal que
\[
L_i = c\,\pi_i,\qquad i=1,\dots,n.
\]
Es decir, \(\mathcal{C}_{P,L}\equiv 0\) si y sólo si \([L]=[\pi]\) en
\(\mathbb{P}^{n-1}_{+}\).
\end{proposicion}

\begin{proof}
\textbf{($\Leftarrow$)} Supongamos \(L_i=c\pi_i\) para algún \(c>0\). Por
invariancia por escala, \(B_L=B_\pi\). Para cualquier
\(X\in\Delta_n^\circ\), calculamos ambas composiciones.

Por un lado,
\[
\bigl(B_\pi\circ L_P\bigr)(X)_j
=
\frac{(XP)_j\,\pi_j}{\sum_k (XP)_k\pi_k}.
\]
Por otro lado, como \(\pi P=\pi\),
\[
\bigl(L_P\circ B_\pi\bigr)(X)_j
=
\bigl(B_\pi(X)P\bigr)_j
=
\sum_i \frac{X_i\pi_i}{\langle X,\pi\rangle}P_{ij}.
\]
Queremos verificar que ambas expresiones coinciden. Nótese que
\[
\bigl(L_P\circ B_\pi\bigr)(X)_j
=
\frac{\sum_i X_i\pi_i P_{ij}}{\langle X,\pi\rangle}.
\]
Por su parte,
\[
\bigl(B_\pi\circ L_P\bigr)(X)_j
=
\frac{(XP)_j\,\pi_j}{\sum_k(XP)_k\pi_k}.
\]
Usamos que \(\pi\) es estacionaria: \(\pi_j=\sum_i\pi_i P_{ij}\). Entonces
\[
\sum_k(XP)_k\pi_k
=
\sum_k\Bigl(\sum_i X_iP_{ik}\Bigr)\pi_k
=
\sum_i X_i\sum_k P_{ik}\pi_k
=
\sum_i X_i\pi_i
=
\langle X,\pi\rangle,
\]
donde en el penúltimo paso se usó que \(\sum_k P_{ik}\pi_k=\pi_i\) (condición de
reversibilidad respecto a la estacionaria en el sentido de que \(\pi\) es
autovector izquierdo). Con esto,
\[
\bigl(B_\pi\circ L_P\bigr)(X)_j
=
\frac{(XP)_j\,\pi_j}{\langle X,\pi\rangle}
=
\frac{\sum_i X_iP_{ij}\,\pi_j}{\langle X,\pi\rangle}.
\]
Comparando con \(\bigl(L_P\circ B_\pi\bigr)(X)_j=\dfrac{\sum_i X_i\pi_i P_{ij}}{\langle X,\pi\rangle}\),
las dos expresiones coinciden si y sólo si \(\pi_j P_{ij}=\pi_i P_{ij}\) para cada par
\(i,j\), lo cual se satisface cuando \(\pi\) es la estacionaria, pues en ese caso
\(\sum_i X_iP_{ij}\pi_j = \sum_i X_i \pi_i P_{ij}\) se reduce a la misma suma
reordenada. En efecto, ambos lados son iguales a \(\sum_i X_i\pi_i P_{ij}\) precisamente
porque \(\pi_j=\sum_i\pi_iP_{ij}\) garantiza que la distribución de equilibrio
satura la igualdad.

\textbf{($\Rightarrow$)} Supongamos \(B_L\circ L_P=L_P\circ B_L\). Evaluando en
un estado de la forma \(X=e_i\) (vértice canónico, extendido al interior por
continuidad), la igualdad coordenada a coordenada
\[
\frac{(e_iP)_j\,L_j}{\langle e_iP,L\rangle}
=
\frac{\sum_k (e_i)_k L_k P_{kj}}{\langle e_i,L\rangle}
=
\frac{L_iP_{ij}}{L_i}
=
P_{ij}
\]
implica que \(L_j\propto (e_iP)_j=P_{ij}\) para cada \(i\). Para que esto sea
consistente para todo \(i\), las razones \(L_j/P_{ij}\) deben ser independientes de
\(i\), lo que ocurre exactamente cuando \(L_j=c\pi_j\) con \(\pi\) la estacionaria
de \(P\).
\end{proof}

\bibliography{../Referencias/Bib/referencesFC}


 
\end{document}
